\section*{Literatura}
\addcontentsline{toc}{section}{Literatura}
\begin{small}\raggedright
\begin{enumerate}[label={},leftmargin=1.2em,itemindent=-1.2em,itemsep=0.15em]
\item Arner P, Bernard S, Salehpour M, Possnert G, Liebl J, Steier P, Buchholz BA, Eriksson M, Arner E, Hauner H, Skurk T, Rydén M, Frayn KN, Spalding KL (2011). Dynamics of human adipose lipid turnover in health and metabolic disease. \emph{Nature} 478(7367):110–113. \href{https://doi.org/10.1038/nature10426}{doi:10.1038/nature10426}
\item Banavar JR, Moses ME, Brown JH, Damuth J, Rinaldo A, Sibly RM, Maritan A (2010). A general basis for quarter-power scaling in animals. \emph{PNAS} 107(36):15816–15820. \href{https://doi.org/10.1073/pnas.1009974107}{doi:10.1073/pnas.1009974107}
\item Bayne BL (1973). Physiological changes in \emph{Mytilus edulis} L. induced by temperature and nutritive stress. \emph{Journal of the Marine Biological Association of the United Kingdom} 53(1):39–58. \href{https://doi.org/10.1017/S0025315400056629}{doi:10.1017/S0025315400056629}
\item Blomberg SP, Garland T, Ives AR (2003). Testing for phylogenetic signal in comparative data: behavioral traits are more labile. \emph{Evolution} 57(4):717–745. \href{https://doi.org/10.1111/j.0014-3820.2003.tb00285.x}{doi:10.1111/j.0014-3820.2003.tb00285.x}
\item Callier V, Nijhout HF (2012). Supply-side constraints are insufficient to explain the ontogenetic scaling of metabolic rate in the tobacco hornworm, \emph{Manduca sexta}. \emph{PLoS ONE} 7(9):e45455. \href{https://doi.org/10.1371/journal.pone.0045455}{doi:10.1371/journal.pone.0045455}
\item Carleton SA, Martínez del Rio C (2005). The effect of cold-induced increased metabolic rate on the rate of $^{13}$C and $^{15}$N incorporation in house sparrows (\emph{Passer domesticus}). \emph{Oecologia} 144(2):226–232. \href{https://doi.org/10.1007/s00442-005-0066-8}{doi:10.1007/s00442-005-0066-8}
\item Forsthoefel DJ, Park AE, Newmark PA (2011). Stem cell-based growth, regeneration, and remodeling of the planarian intestine. \emph{Developmental Biology} 356(2):445–459. \href{https://doi.org/10.1016/j.ydbio.2011.05.669}{doi:10.1016/j.ydbio.2011.05.669}
\item Gabbott PA, Bayne BL (1973). Biochemical effects of temperature and nutritive stress on \emph{Mytilus edulis} L. \emph{Journal of the Marine Biological Association of the United Kingdom} 53(2):269–286. \href{https://doi.org/10.1017/S0025315400022268}{doi:10.1017/S0025315400022268}
\item Gazit Y, Berk DA, Leunig M, Baxter LT, Jain RK (1995). Scale-invariant behavior and vascular network formation in normal and tumor tissue. \emph{Physical Review Letters} 75(12):2428–2431. \href{https://doi.org/10.1103/PhysRevLett.75.2428}{doi:10.1103/PhysRevLett.75.2428}
\item Gehr P, Mwangi DK, Ammann A, Maloiy GMO, Taylor CR, Weibel ER (1981). Design of the mammalian respiratory system. V. Scaling morphometric pulmonary diffusing capacity to body mass: wild and domestic mammals. \emph{Respiration Physiology} 44(1):61–86. \href{https://doi.org/10.1016/0034-5687(81)90077-3}{doi:10.1016/0034-5687(81)90077-3}
\item Glazier DS (2010). A unifying explanation for diverse metabolic scaling in animals and plants. \emph{Biological Reviews} 85(1):111–138. \href{https://doi.org/10.1111/j.1469-185X.2009.00095.x}{doi:10.1111/j.1469-185X.2009.00095.x}
\item Glazier DS, Hirst AG, Atkinson D (2015). Shape shifting predicts ontogenetic changes in metabolic scaling in diverse aquatic invertebrates. \emph{Proc R Soc B} 282:20142302. \href{https://doi.org/10.1098/rspb.2014.2302}{doi:10.1098/rspb.2014.2302}
\item Haller M, Brechtelsbauer H, Finsterer U, Forst H, Bein T, Briegel J, Peter K (1992). Bestimmung des Plasmavolumens mit Indocyaningrün beim Menschen [The determination of plasma volume using indocyanine green in man]. \emph{Der Anaesthesist} 41(3):115–120. \href{https://pubmed.ncbi.nlm.nih.gov/1570882/}{PMID:1570882}
\item Heiling VJ, Miles JM, Jensen MD (1991). How valid are isotopic measurements of fatty acid oxidation? \emph{American Journal of Physiology – Endocrinology and Metabolism} 261(5):E572–E577. \href{https://doi.org/10.1152/ajpendo.1991.261.5.E572}{doi:10.1152/ajpendo.1991.261.5.E572}
\item Helm BR, Davidowitz G (2013). Mass and volume growth of an insect tracheal system within a single instar. \emph{Journal of Experimental Biology} 216(24):4703–4711. \href{https://doi.org/10.1242/jeb.080648}{doi:10.1242/jeb.080648}
\item Henderson D (1977). The effect of worm, age, weight and number in the infection on the absorption of glucose by \emph{Hymenolepis diminuta}. \emph{Parasitology} 75(3):277–284. \href{https://doi.org/10.1017/S0031182000051817}{doi:10.1017/S0031182000051817}
\item Hirst AG, Glazier DS, Atkinson D (2014). Body shape shifting during growth permits tests that distinguish between competing geometric theories of metabolic scaling. \emph{Ecology Letters} 17(10):1274–1281. \href{https://doi.org/10.1111/ele.12334}{doi:10.1111/ele.12334}
\item Hoppeler H, Mathieu O, Krauer R, Claassen H, Armstrong RB, Weibel ER (1981a). Design of the mammalian respiratory system. VI. Distribution of mitochondria and capillaries in various muscles. \emph{Respiration Physiology} 44(1):87–111. \href{https://doi.org/10.1016/0034-5687(81)90078-5}{doi:10.1016/0034-5687(81)90078-5}
\item Hoppeler H, Mathieu O, Weibel ER, Krauer R, Lindstedt SL, Taylor CR (1981b). Design of the mammalian respiratory system. VIII. Capillaries in skeletal muscles. \emph{Respiration Physiology} 44(1):129–150. \href{https://doi.org/10.1016/0034-5687(81)90080-3}{doi:10.1016/0034-5687(81)90080-3}
\item Kearney MR, White CR (2012). Testing metabolic theories. \emph{American Naturalist} 180(5):546–565. \href{https://doi.org/10.1086/667860}{doi:10.1086/667860}
\item Kelly TL, Wilson KE, Heymsfield SB (2009). Dual energy X-ray absorptiometry body composition reference values from NHANES. \emph{PLoS ONE} 4(9):e7038. \href{https://doi.org/10.1371/journal.pone.0007038}{doi:10.1371/journal.pone.0007038}
\item Kooijman SALM (1986). Energy budgets can explain body size relations. \emph{Journal of Theoretical Biology} 121(3):269–282. \href{https://doi.org/10.1016/S0022-5193(86)80107-2}{doi:10.1016/S0022-5193(86)80107-2}
\item Kooijman SALM (2014). Metabolic acceleration in animal ontogeny: an evolutionary perspective. \emph{Journal of Sea Research} 94:128–137. \href{https://doi.org/10.1016/j.seares.2014.06.005}{doi:10.1016/j.seares.2014.06.005}
\item Lefevre S, McKenzie DJ, Nilsson GE (2017). Models projecting the fate of fish populations under climate change need to be based on valid physiological mechanisms. \emph{Global Change Biology} 23(9):3449–3459. \href{https://doi.org/10.1111/gcb.13652}{doi:10.1111/gcb.13652}
\item Lorthois S, Cassot F (2010). Fractal analysis of vascular networks: insights from morphogenesis. \emph{Journal of Theoretical Biology} 262(4):614–633. \href{https://doi.org/10.1016/j.jtbi.2009.10.037}{doi:10.1016/j.jtbi.2009.10.037}
\item Lupatsch I, Kissil GWm, Sklan D (2003). Defining energy and protein requirements of gilthead seabream (\emph{Sparus aurata}) to optimize feeds and feeding regimes. \emph{Israeli Journal of Aquaculture – Bamidgeh} 55(4):243–257. \href{https://doi.org/10.46989/001c.20354}{doi:10.46989/001c.20354}
\item Maino JL, Kearney MR (2014). Ontogenetic and interspecific metabolic scaling in insects. \emph{American Naturalist} 184(6):695–701. \href{https://doi.org/10.1086/678401}{doi:10.1086/678401}
\item Maino JL, Kearney MR (2015a). Testing mechanistic models of growth in insects. \emph{Proc R Soc B} 282:20151973. \href{https://doi.org/10.1098/rspb.2015.1973}{doi:10.1098/rspb.2015.1973}
\item Maino JL, Kearney MR (2015b). Ontogenetic and interspecific scaling of consumption in insects. \emph{Oikos} 124(12):1564–1570. \href{https://doi.org/10.1111/oik.02341}{doi:10.1111/oik.02341}
\item Maino JL, Kearney MR, Nisbet RM, Kooijman SALM (2014). Reconciling theories for metabolic scaling. \emph{Journal of Animal Ecology} 83(1):20–29. \href{https://doi.org/10.1111/1365-2656.12085}{doi:10.1111/1365-2656.12085}
\item Mainster MA (1990). The fractal properties of retinal vessels: embryological and clinical implications. \emph{Eye} 4(1):235–241. \href{https://doi.org/10.1038/eye.1990.33}{doi:10.1038/eye.1990.33}
\item Marques GM, Augustine S, Lika K, Pecquerie L, Domingos T, Kooijman SALM (2018). The AmP project: comparing species on the basis of dynamic energy budget parameters. \emph{PLoS Computational Biology} 14(5):e1006100. \href{https://doi.org/10.1371/journal.pcbi.1006100}{doi:10.1371/journal.pcbi.1006100}
\item Masters BR (2004). Fractal analysis of the vascular tree in the human retina. \emph{Annual Review of Biomedical Engineering} 6:427–452. \href{https://doi.org/10.1146/annurev.bioeng.6.040803.140100}{doi:10.1146/annurev.bioeng.6.040803.140100}
\item Monahan-Earley R, Dvorak AM, Aird WC (2013). Evolutionary origins of the blood vascular system and endothelium. \emph{Journal of Thrombosis and Haemostasis} 11(Suppl 1):46–66. \href{https://doi.org/10.1111/jth.12253}{doi:10.1111/jth.12253}
\item Oikawa S, Itazawa Y (1985). Gill and body surface areas of the carp in relation to body mass, with special reference to the metabolism-size relationship. \emph{Journal of Experimental Biology} 117:1–14. \href{https://doi.org/10.1242/jeb.117.1.1}{doi:10.1242/jeb.117.1.1}
\item Ojeda-Avila T, Woods HA, Raguso RA (2003). Effects of dietary variation on growth, composition, and maturation of \emph{Manduca sexta} (Sphingidae: Lepidoptera). \emph{Journal of Insect Physiology} 49(4):293–306. \href{https://doi.org/10.1016/S0022-1910(03)00003-9}{doi:10.1016/S0022-1910(03)00003-9}
\item Okie JG (2013). General models for the spectra of surface area scaling strategies of cells and organisms: fractality, geometric dissimilitude, and internalization. \emph{American Naturalist} 181(3):421–439. \href{https://doi.org/10.1086/669150}{doi:10.1086/669150}
\item Pagel M (1999). Inferring the historical patterns of biological evolution. \emph{Nature} 401:877–884. \href{https://doi.org/10.1038/44766}{doi:10.1038/44766}
\item Palzenberger M, Pohla H (1992). Gill surface area of water-breathing freshwater fish. \emph{Reviews in Fish Biology and Fisheries} 2(3):187–216. \href{https://doi.org/10.1007/BF00045037}{doi:10.1007/BF00045037}
\item Park Y, Grellner WJ, Harris WS, Miles JM (2000). A new method for the study of chylomicron kinetics in vivo. \emph{American Journal of Physiology – Endocrinology and Metabolism} 279(6):E1258–E1263. \href{https://doi.org/10.1152/ajpendo.2000.279.6.E1258}{doi:10.1152/ajpendo.2000.279.6.E1258}
\item Pauly D (2021). The gill-oxygen limitation theory (GOLT) and its critics. \emph{Science Advances} 7(2):eabc6050. \href{https://doi.org/10.1126/sciadv.abc6050}{doi:10.1126/sciadv.abc6050}
\item Pauly D, Cheung WWL (2018). Sound physiological knowledge and principles in modeling shrinking of fishes under climate change. \emph{Global Change Biology} 24(1):e15–e26. \href{https://doi.org/10.1111/gcb.13831}{doi:10.1111/gcb.13831}
\item Reed DR, Bachmanov AA, Tordoff MG (2007). Forty mouse strain survey of body composition. \emph{Physiology \& Behavior} 91(5):593–600. \href{https://doi.org/10.1016/j.physbeh.2007.03.026}{doi:10.1016/j.physbeh.2007.03.026}
\item Riisgård HU, Larsen PS, Pleissner D (2014). Allometric equations for maximum filtration rate in blue mussels \emph{Mytilus edulis} and importance of condition index. \emph{Helgoland Marine Research} 68(1):193–198. \href{https://doi.org/10.1007/s10152-013-0377-9}{doi:10.1007/s10152-013-0377-9}
\item Rombough PJ, Moroz BM (1990). The scaling and potential importance of cutaneous and branchial surfaces in respiratory gas exchange in young chinook salmon (\emph{Oncorhynchus tshawytscha}). \emph{Journal of Experimental Biology} 154:1–12. \href{https://doi.org/10.1242/jeb.154.1.1}{doi:10.1242/jeb.154.1.1}
\item Savage VM, Deeds EJ, Fontana W (2008). Sizing up allometric scaling theory. \emph{PLoS Computational Biology} 4(9):e1000171. \href{https://doi.org/10.1371/journal.pcbi.1000171}{doi:10.1371/journal.pcbi.1000171}
\item Taylor CR, Weibel ER (1981). Design of the mammalian respiratory system. I. Problem and strategy. \emph{Respiration Physiology} 44(1):1–10. \href{https://doi.org/10.1016/0034-5687(81)90073-6}{doi:10.1016/0034-5687(81)90073-6}
\item Tessier AJ, Henry LL, Goulden CE, Durand MW (1983). Starvation in \emph{Daphnia}: energy reserves and reproductive allocation. \emph{Limnology and Oceanography} 28(4):667–676. \href{https://doi.org/10.4319/lo.1983.28.4.0667}{doi:10.4319/lo.1983.28.4.0667}
\item Thomas SM, Crowther TW (2015). Predicting rates of isotopic turnover across the animal kingdom: a synthesis of existing data. \emph{Journal of Animal Ecology} 84(3):861–870. \href{https://doi.org/10.1111/1365-2656.12326}{doi:10.1111/1365-2656.12326}
\item Thommen A, Werner S, Frank O, Philipp J, Knittelfelder O, Quek Y, Fahmy K, Shevchenko A, Friedrich BM, Jülicher F, Rink JC (2019). Body size-dependent energy storage causes Kleiber's law scaling of the metabolic rate in planarians. \emph{eLife} 8:e38187. \href{https://doi.org/10.7554/eLife.38187}{doi:10.7554/eLife.38187}
\item Threlkeld ST (1976). Starvation and the size structure of zooplankton communities. \emph{Freshwater Biology} 6(6):489–496. \href{https://doi.org/10.1111/j.1365-2427.1976.tb01640.x}{doi:10.1111/j.1365-2427.1976.tb01640.x}
\item Weibel ER, Hoppeler H (2005). Exercise-induced maximal metabolic rate scales with muscle aerobic capacity. \emph{Journal of Experimental Biology} 208(9):1635–1644. \href{https://doi.org/10.1242/jeb.01548}{doi:10.1242/jeb.01548}
\item Weibel ER, Taylor CR, Hoppeler H (1991). The concept of symmorphosis: a testable hypothesis of structure-function relationship. \emph{PNAS} 88(22):10357–10361. \href{https://doi.org/10.1073/pnas.88.22.10357}{doi:10.1073/pnas.88.22.10357}
\item West GB, Brown JH, Enquist BJ (1997). A general model for the origin of allometric scaling laws in biology. \emph{Science} 276(5309):122–126. \href{https://doi.org/10.1126/science.276.5309.122}{doi:10.1126/science.276.5309.122}
\item West GB, Brown JH, Enquist BJ (1999). The fourth dimension of life: fractal geometry and allometric scaling of organisms. \emph{Science} 284(5420):1677–1679. \href{https://doi.org/10.1126/science.284.5420.1677}{doi:10.1126/science.284.5420.1677}
\item White CR, Kearney MR, Matthews PGD, Kooijman SALM, Marshall DJ (2011). A manipulative test of competing theories for metabolic scaling. \emph{American Naturalist} 178(6):746–754. \href{https://doi.org/10.1086/662666}{doi:10.1086/662666}
\end{enumerate}
\end{small}
